\subsection{Contamination Detection Methods}
\label{sec:methods}

We employ ten contamination detection methods, each producing a scalar score per code sample that is thresholded into a binary contamination decision. The methods differ in the statistical signal they exploit and in whether they require only the code (\texttt{Response}) or also the natural-language specification (\texttt{Instruction}). Below we describe each method in turn.

\paragraph{Zlib.}
Zlib exploits the dual signal that training-seen code tends to be both low-perplexity and highly compressible. For each sample the target model computes the mean token log-probability $\ell$ (a negative quantity, closer to zero for familiar code), and the code bytes are zlib-compressed to length $c$. The score is defined as $\mathrm{zlib}=\ell/c$. Memorized code exhibits a \emph{larger} ratio because it is simultaneously likely (high $\ell$) and redundant (small $c$). A sample is classified as contaminated when $\mathrm{zlib}\ge -0.003$.

\paragraph{Min-K\% ($\mathrm{mink}_{0.2}$).}
Min-K\% Prob relies on the observation that in memorized sequences even the least likely tokens receive abnormally high probabilities. The token log-probabilities are sorted in ascending order, and the mean of the lowest $20\%$ is taken as the score. If the model has seen the code, the tail tokens are also lifted, yielding a higher score. A sample is contaminated when $\mathrm{mink}_{0.2}\ge -3.0$.

\paragraph{Min-K\%++ ($\mathrm{mink}^{\!++}_{0.2}$).}
Min-K\%++ refines Min-K\% by normalizing the token log-probabilities with the per-sequence mean and variance before selecting the lowest $20\%$. This standardization removes scale differences across sequences of different lengths and styles, making the score more comparable. A sample is contaminated when $\mathrm{mink}^{\!++}_{0.2}\ge -3.0$.

\paragraph{DC-PDD.}
DC-PDD calibrates the target model's likelihood against a reference corpus so as to remove the confound of ``commonness''. A smoothed token-frequency distribution is built from an external Verilog corpus (\texttt{ref\_verilog.jsonl}), and a cross-entropy-style divergence between the sample's token probabilities and this reference is computed and negated. A sample whose model likelihood exceeds what its generic frequency warrants is flagged. A sample is contaminated when $\mathrm{DCPDD}\le -3.0$.

\paragraph{Perplexity ($\mathrm{ppl}$).}
PPL uses the target model's likelihood of the code as the signal. Although named after perplexity, the score is in fact the mean token log-probability $\bar\ell$ (a negative quantity) rather than the true perplexity $\exp(-\bar\ell)$; the two are monotonically inversely related. Memorized code has higher mean log-probability (equivalently lower perplexity), hence a higher score. A sample is contaminated when $\mathrm{ppl}\ge -2.2$.

\paragraph{PETAL.}
PETAL is a black-box-friendly method that calibrates the natural effect of semantic similarity on log-probability. A linear mapping from semantic similarity to token log-probability is first fitted on a surrogate model; the target model's within-code span similarities are then mapped to estimated log-probabilities, and the score is $\mathrm{PETAL}=-\overline{\hat\ell}$. Because memorization raises likelihood beyond what semantic similarity alone explains, this score detects memory. A sample is contaminated when $\mathrm{PETAL}<0.8$. PETAL uses a surrogate model (Qwen2.5-1.5B-Instruct) and an embedding model (all-MiniLM-L6-v2).

\paragraph{Neighbor.}
Neighbor compares the likelihood of the target code against its near-duplicate variants. Neighbors are generated by CodeBERT masked-language modeling with consistent identifier renaming, and the score is $\mathrm{neighbor}=\ell_{\mathrm{main}}-\overline{\ell_{\mathrm{neigh}}}$. Memorized code stands out sharply against its perturbed versions. A sample is contaminated when $\mathrm{neighbor}\ge 0.8$.

\paragraph{Reference.}
Reference calibrates the target model's likelihood against a reference model. The rationale is that the reference model is unlikely to have seen the specific code samples under test, so its likelihood reflects an un-memorized baseline; a target-model likelihood substantially higher than this baseline indicates extra memorization of that code. The score is $\mathrm{ref}=\ell_{\mathrm{main}}-\overline{\ell_{\mathrm{ref}}}$. A sample is contaminated when $\mathrm{ref}\ge 0.13$. The reference model is Qwen2.5-1.5B-Instruct.

\paragraph{Token Overlap.}
Token Overlap tests whether the model can reproduce the true suffix from the prefix. The code is split into a prefix (first piece) and a true suffix (second piece); the target model then generates a completion from the prefix, where the guided prompt additionally prepends the sample's Instruction (mimicking the ``instruction $\rightarrow$ code'' task seen in training) and the general prompt uses only the prefix as an instruction-free control. ROUGE-L measures the token overlap of each completion against the true suffix, with the split ratio $\mathrm{min\_p}=\mathrm{max\_p}=20$ (the true suffix is the last $20\%$) and $\mathrm{max\_new\_tokens}=128$. The intuition is that if the model has seen the code during training---especially together with its Instruction---it can faithfully ``recall'' the true suffix given the prefix, yielding markedly higher overlap. The guided-completion similarity is used as the final score for thresholding; the general similarity serves mainly as a statistical-test control (testing whether guided significantly exceeds general, i.e., whether the Instruction genuinely aids recall) and is not directly used in downstream decisions. A sample is contaminated when the score $\ge 0.55$.

\paragraph{CDD.}
CDD detects memorization by measuring the determinism of the model's generation. The core intuition is that a memorized sample (whose ``instruction $\rightarrow$ code'' mapping was seen in training) is generated with high determinism, so even stochastic sampling produces nearly identical code, whereas an unmemorized sample is generated with uncertainty and sampling yields diverse outputs. Given the sample's Instruction, the target model performs one greedy decoding (a single deterministic output) and $\mathrm{num\_samples}=50$ sampled decodings (temperature $0.8$, top-$p$ $1.0$, $\mathrm{max\_new\_tokens}=128$). For each sampled output, its edit distance to the greedy output is computed; a distance below $\alpha\times\mathrm{max\_len}$ ($\alpha=0.05$, where $\mathrm{max\_len}$ is the token length of the longer of the greedy and sampled outputs, i.e., $5\%$ of the code length) is treated as agreement. The contamination score is the fraction of sampled outputs that agree with the greedy output, approaching 1 for memorized code and 0 otherwise. A sample is contaminated when the score exceeds $\xi=0.01$. CDD consumes both \texttt{Instruction} and \texttt{Response}.

\begin{table}[t]
\centering
\caption{Summary of the ten detection methods. Input indicates whether the method uses the code (\texttt{Response}), the instruction (\texttt{Instruction}), or auxiliary models/corpora.}
\label{tab:methods}
\small
\begin{tabular}{llcl}
\toprule
Method & Input & Direction & Threshold \\
\midrule
Zlib & Response & $\ge$ & $-0.003$ \\
Min-K\% & Response & $\ge$ & $-3.0$ \\
Min-K\%++ & Response & $\ge$ & $-3.0$ \\
DC-PDD & Response + reference corpus & $\le$ & $-3.0$ \\
Perplexity & Response & $\ge$ & $-2.2$ \\
PETAL & Response + surrogate/embedding models & $<$ & $0.8$ \\
Neighbor & Response + CodeBERT & $\ge$ & $0.8$ \\
Reference & Response + reference model & $\ge$ & $0.13$ \\
Token Overlap & Response (guided part uses Instruction) & $\ge$ & $0.55$ \\
CDD & Instruction + Response & $>$ & $0.01$ \\
\bottomrule
\end{tabular}
\end{table}

Note that only CDD explicitly uses the \texttt{Instruction}; Token Overlap's guided completion also depends on the \texttt{Instruction}, whereas the remaining eight methods operate solely on the \texttt{Response} (code). This input dependency implies that when the negative-sample instructions are regenerated, only the CDD and Token Overlap features need to be recomputed.
